% 1. Hafta 
% -> Neden Python?
% -> Program Introduction
% -> IDE kurulumu
% -> Programlama ve Pythona giriş 
% 	- Variables
% 	- Console input output

% pdflatex --shell-escape lecture1.tex & bibtex lecture1 & pdflatex --shell-escape lecture1.tex & pdflatex --shell-escape lecture1.tex

\input{../KU-Beamer-Template/style/template.tex}
\usepackage{../KU-Beamer-Template/style/koc} 
\usepackage{minted}
\usepackage{upquote}
\usepackage{animate}

\title{KOLT Python} 
\subtitle{Introduction} 
\newdate{date}{23}{09}{2019}
\date{\displaydate{date}}
\author{Ahmet Uysal}

\titlegraphic{\includegraphics[scale=0.2]{../KU-Beamer-Template/style/images/logo_kolt.pdf}}

\setbeamercovered{invisible} % transparent

\begin{document}
  \maketitle

	\frame{\frametitle{Agenda}\tableofcontents}

	\section{Program Information}

	\begin{frame}{Course Outcomes}
		\LARGE
		\begin{itemize}
			\pause
			\item Apply basic programming concepts using Python
			\pause
			\item Demonstrate how Python can be used in different areas or disciplines
			\pause
			\item Create code that is easy to understand
			\pause
			\item \textbf{Implement practical challenges} by gaining experience in Python
		\end{itemize}
	\end{frame}

	\begin{frame}{Why Python?}
		\pause
		\begin{itemize}
			\LARGE
			\item Easy Syntax
			\pause
			\item Beginner Friendly
			-most popular language for introductory CS courses in top universities\cite{guo1}- 
			\pause
			\item Wide usage area
			\pause
			\item Large and growing community
		\end{itemize}
	\end{frame}

	\begin{frame}{Some of the Usage Areas \cite{survey_jetbrains2018}}
		\begin{itemize}
			\pause
			\LARGE
			\item Data Analysis
			\pause
			\item Web Development
			\pause
			\item System Administration
			\pause
			\item Machine Learning
			\pause
			\item Web Parsers/Scrawlers
			\pause
			\item Testing
			\pause
			\item Education
			\pause
			\item Network Programming
			\pause
			\item ...
		\end{itemize}	
	\end{frame}
	
	\begin{frame}{Python at Koç University}
		\begin{itemize}
			\LARGE
			\visible<2->{
			\item COMP341: Introduction to Artificial Intelligence
			}
			\visible<3->{
			\item ENGR421/521: Introduction to Machine Learning
			}
			\visible<4->{
			\item ENGR350 (Selected Topics - Summer18/Spring19): Introduction to Programming for Data Science
			}
			\visible<5->{
			\item INTL450 (Selected Topics - Spring19): Advanced Data Analysis in Python
			}
		\end{itemize}
		\begin{columns}
			\column{0.3\textwidth}
			\visible<2->{
				\includegraphics[width=\textwidth]{images/berkeley_pacman.png}
			}
			\column{0.3\textwidth}
			\visible<3->{
				\includegraphics[width=\textwidth]{images/nn.jpg}
			}
			\column{0.3\textwidth}
			\visible<5->{
				\includegraphics[width=\textwidth]{images/analysis.jpeg}
			}
		\end{columns}
	\end{frame}

	\begin{frame}{Python at Industry}
		\begin{table}[]
			\resizebox{\textwidth}{!}{
			\begin{tabular}{lll}
				\includegraphics[width=2cm]{images/firefox.png}& \includegraphics[width=3cm]{images/google.png} & \includegraphics[width=2cm]{images/nasa.png} \\
				\includegraphics[width=2cm]{images/instagram.png}& \includegraphics[width=3cm]{images/youtube.png} & \includegraphics[width=2cm]{images/reddit.png}\\
			\end{tabular}}
		\end{table}
	\end{frame}

	\begin{frame}{Python at Industry \cite{instagram_engineering}}
		\centering
		\begin{tabular}{c}
		\includegraphics[width=0.5\textwidth]{images/instagram_python.png}\\
		\includegraphics[width=0.6\textwidth]{images/nasa_python.png}
		\end{tabular}
	\end{frame}

    \begin{frame}{Python Everywhere}
		\centering
		\includegraphics[width=0.8\textwidth]{images/python_everywhere.jpg}
	\end{frame}


	\section{Logistics}
		\begin{frame}{Who Are We?}
			\begin{columns}
				\column{5cm}
				\centering 
				\visible<1->{
					\includegraphics[width=2.0cm]{images/ahmet.png}\\
					Ahmet Uysal\\
					auysal16@ku.edu.tr\\
				}
				\visible<3->{
					\includegraphics[width=2.0cm]{images/gulsena.jpg}\\
					Gül Sena Altıntaş\\
					galtintas17@ku.edu.tr\\
				}
				\column{5cm}
				\centering
				\visible<2->{
					\includegraphics[width=2.0cm]{images/ceren.jpg}\\
					Ceren Kocaoğullar\\
					ckocaogullar15@ku.edu.tr
				}
				\visible<4->{
					\includegraphics[width=2.0cm]{images/hasancan.jpg}\\
					Hasan Can Aslan\\
					haslan16@ku.edu.tr
				}
			\end{columns}
		\end{frame}

		\begin{frame}{What Will We Do?}
			\LARGE
			\pause
			\texttt{\textbf{Lecture} Monday 11:30-12:45}\\
			\pause
			\texttt{\textbf{Section} Wednesday 14:30-15:45
			\textbf{or}}\\
			\texttt{\textbf{Section} Thursday 11:30-12:45
			\textbf{or}}\\
			\texttt{\textbf{Section} Thursday 13:00-14:15
			\textbf{or}}\\
			\texttt{\textbf{Section} Thursday 16:00-17:15}\\
		\end{frame}

		\begin{frame}{Programming Assignments(Don't Be Afraid!)}
			\pause
			\begin{itemize}
				\LARGE
				\item 4-6 in total
				\pause
				\item Review sessions to \textbf{help you}
				\pause
				\item Some assignments will have \textbf{autograders} to help you find your mistakes and test your code.
				\pause
				\item Later assignments will be based on \textbf{your interests!}
				\pause
				%\item We can also organize hackathons if requested :)
			\end{itemize}
			\ifthenelse{\boolean{beamer@anotherslide}}{
				\phantom{\rule{0.3\linewidth}{0.3\linewidth}}
			}
			{
			
				\begin{center}
					\animategraphics[loop,autoplay,width=0.3\linewidth]{12}{gifs/assignment/}{0}{239}				
				\end{center}
			}
			
		\end{frame}

		\begin{frame}{Certificate Requirements}
			\pause
			\begin{itemize}
				\LARGE
				\item At most \textbf{3 unexcused absences}, including sections.
				\pause
				\item Working on and submitting all homework assignments. Submissions that do not pass the autograders will be examined by us.
				\pause
				\item We do not expect that you ace all programming assignments. But, we expect that you \textbf{spend time} on them!
				\pause
				\item Complying to \href{https://vpaa.ku.edu.tr/academic/student-code-of-conduct}{\underline{\textit{Koç University Code of Conduct}}}. 
			\end{itemize}
		\end{frame}

	\section{Installations}

	 	\begin{frame}{Installations}
			\LARGE
			\pause
			NO INSTALLATIONS!
			\begin{center}
				\includegraphics[width=3.0cm]{images/gotcha_jesus.jpg}
			\end{center}
			\pause
			We will use \href{https://repl.it}{Repl.it}, an online IDE.
			\begin{center}
				\includegraphics[width=3.0cm]{images/replit_logo.png}
			\end{center}
		\end{frame}
		
		\begin{frame}{Sign Up to Repl.it} 		
			\LARGE
			\pause
			We already sent you an invite, check your email.
			\pause
			\begin{center}
				\includegraphics[width=\textwidth]{images/replit_email.png}
			\end{center}
		\end{frame}

		\begin{frame}{Sign Up Flow}
		
		\begin{columns}
			\pause
			\begin{column}{0.33\textwidth}
				\centering
				\includegraphics[width=\textwidth]{images/replit_gauth1.png}
			\end{column}
			\pause
			\begin{column}{0.33\textwidth}
				\centering
				\includegraphics[width=\textwidth]{images/replit_gauth2.png}
			\end{column}
			\pause
			\begin{column}{0.33\textwidth}
				\centering
				\includegraphics[width=\textwidth]{images/replit_success.png}
			\end{column}
		\end{columns}	
		\end{frame}

		\begin{frame}{Repl.it}
				\begin{tikzpicture}
					\node[anchor=north west,inner sep=0] at (0,0) {\includegraphics[width=1\textwidth]{images/replit_interface.png}};
					\draw<2>[red,ultra thick,rounded corners] (0.5\textwidth,0) rectangle (\textwidth,-4.25);				
					\draw<3>[red,ultra thick,rounded corners] (0,0) rectangle (0.5\textwidth,-2.9);
					\draw<4>[red,ultra thick,rounded corners] (2.2,0) rectangle (3.2,-0.3);
					\draw<5>[red,ultra thick,rounded corners] (0,-2.9) rectangle (0.5\textwidth,-4.25);
				\end{tikzpicture}
		\end{frame}

	\section{Introduction}
		
		\begin{frame}{Comments}
			\LARGE
			\inputminted[frame=single,framesep=2pt]{python3}{code-examples/comments.py}
			Python will basically ignore comments, they are purely written \textbf{for humans}!
		\end{frame}

		\begin{frame}{Variables}
			\LARGE
			\begin{itemize}
				\pause
				\item How to represent/store values in Python?
				\pause
				\item The answer is:
				\item \huge{\textbf{VARIABLES!}}
				\vspace{5mm}
				\pause
				\item But what the heck are variables?
			\end{itemize}
		\end{frame}
		
		\begin{frame}{Variables}
			\begin{columns}[T]
				\begin{column}{0.2\textwidth}
					\centering
					Think of them as
					\pause					
				\end{column}
				\begin{column}{0.8\textwidth}
					\centering
					\includegraphics[width=\textwidth]{images/box.png}
				\end{column}
			\end{columns}
		\end{frame}
		
		\begin{frame}{Variables}
			\begin{columns}[T]
				\begin{column}{0.6\textwidth}
					\centering
					Tiny boxes in the computer memory
					\pause					
				\end{column}
				\begin{column}{0.4\textwidth}
					\centering
					\includegraphics[width=\textwidth]{images/memory.jpg}
				\end{column}
			\end{columns}
		\end{frame}
		
		\begin{frame}{Variables}
			\LARGE
			\begin{itemize}
				\item Which kind of values we need to keep in these boxes (variables)?
				\pause
					\begin{itemize}
						\Large
						\item Numbers?
						\pause
						\item Texts?
						\pause
						\item Individual Characters?
						\pause
						\item Starting time of the class?
						\pause
						\item Colors?
						\pause
						\item Truth Values?
						\pause
						\item People?
					\end{itemize} 
			\end{itemize}
		\end{frame}

		\begin{frame}{Variables}
			\LARGE
			\begin{table}[]
				\resizebox{\textwidth}{!}{
				\begin{tabular}{|l|l|l|}
				\hline
				\textcolor{koc}{\textbf{Type}} &\textcolor{koc}{\textbf{Explanation}} & \textcolor{koc}{\textbf{Examples}} \\ \hline
				\textbf{\texttt{int}}  & represent \textbf{integers} & 3, 4, 17, -10 \\ \hline 
				\textbf{\texttt{float}} & represent \textbf{real numbers} & 3.0, 1.11, -109.123123 \\ \hline
				\textbf{\texttt{bool}} & represent \textbf{boolean} truth values & \texttt{True}, \texttt{False} \\ \hline
				\textbf{\texttt{str}} & A sequence of characters. & \textquotesingle Hello\textquotesingle, \textquotesingle \textquotesingle, \textquotesingle 3\textquotesingle \\ \hline
				\textbf{\texttt{NoneType}}& special and has one value, None & \texttt{None} \\ \hline
				\end{tabular}}
			\end{table}
			\pause
			OK, but how do we create one?
		\end{frame}

		\begin{frame}{Variables}
			\inputminted[frame=single,framesep=2pt,fontsize=\LARGE]{python3}{code-examples/variables.py}
		\end{frame}
		
		
		\begin{frame}{Variables}
		    \begin{itemize}
		        \item \LARGE BE CAREFUL!
		        \vspace{2mm}
				\pause
				\item \large
				Andy could fit all his toys into one box, but this is not the Toy Story.
				\pause
				\item \LARGE
				These variable boxes can keep only \textbf{one thing at a time.}
		    \end{itemize}
			\vspace{2mm}
		    \centering
            \includegraphics[width=0.6\textwidth]{images/box_many.jpg}
		\end{frame}
		
		\begin{frame}{How about types of variables?}
			\LARGE
			Special method called \texttt{\textbf{type()}} 
			\inputminted[frame=single,framesep=2pt]{python3}{code-examples/types.py}
			Python knows variables' type even if you don't know it!
		\end{frame}

		\begin{frame}{Console I/O(Input/Output)}
			\LARGE
			Now we can store the data we know, \\
			\pause
			how about interacting with the user? \\
			\pause
			\texttt{print(), input()}
			\pause
			\inputminted[frame=single,framesep=2pt]{python3}{code-examples/io.py}
		\end{frame}

		\begin{frame}{Console I/O(Input/Output)}
			\huge
			\textbf{\texttt{print(*args, sep=\textquotesingle \ \textquotesingle, end=\textquotesingle \textbackslash n\textquotesingle )}}
			\pause
			\begin{itemize}
				\LARGE
				\item Can take arbitrary number of arguments
				\pause
				\item Separates elements with space by default
				\pause
				\item Adds newline character \texttt{\textquotesingle \textbackslash n\textquotesingle} to end by default
			\end{itemize}
			
			\pause
			\textbf{\texttt{input([prompt])}}
			\pause
			\begin{itemize}
				\LARGE
				\item Prints the prompt to Console
				\pause
				\item Program is paused until user enters something
				\pause
				\item \textbf{returns an \texttt{str} object!} 
			\end{itemize}
		\end{frame}

		\begin{frame}{Example Program}
			\inputminted[frame=single,framesep=2pt,fontsize=\LARGE]{python3}{code-examples/example_io.py}
			\pause
			\centering
			\includegraphics[width=0.25\textwidth]{images/thanks.png}

		\end{frame}
		
		\begin{frame}{Get Your Hands Dirty}
		    \vspace{5mm}
		    \centering
            \includegraphics[width=0.8\textwidth]{images/your_turn.jpg}

		\end{frame}

	\section{References}
		\begin{frame}{References}
			\bibliographystyle{IEEETran}
			\bibliography{references}
		\end{frame}
\end{document}